% !PRINTED 470
\weberpage{470}

%<<CONT
or, more precisely, according to the choice of sign, into two equations which, apart from the product
\[
\sqrt{2}\, f(\omega)^2
\]
contain only rational numerical coefficients. One sees the latter from the fact that in equation (7) there occur only such products $u^h v^k$ in which $h+k$ is even (cf. the beginning of § 74).

Relation (14) now requires that $\omega$ satisfy an equation
\weq{16}{\frac{c + \partial\omega}{a} = \frac{\gamma(\omega+1) + \delta(\omega-1)}{\alpha(\omega+1) + \beta(\omega-1)}}
in which $\begin{pmatrix} \alpha, & \beta \\ \gamma, & \delta \end{pmatrix}$ is a linear transformation belonging to the first or the second class. But in order that one of the two equations (14) be actually satisfied, there is further required, by § 40, (12):
\weq{17}{(\alpha - \beta)(\alpha + \beta + \gamma - \delta) \equiv 0 \quad (\mathrm{mod}\ 8).}
Now if $\omega$ is a root of the quadratic equation (2), it follows from (16), as above:
\weq{18}{
\begin{aligned}
\partial(\alpha + \beta) &= Ax, \\
\partial(\alpha - \beta) &= Bx + y, \\
c(\alpha - \beta) - a(\gamma - \delta) &= Cx, \\
c(\alpha + \beta) - a(\gamma + \delta) &= Bx - y, \\
2n &= y^2 + mx^2.
\end{aligned}
}
If now, as was assumed,
\[
m \equiv 1 \quad (\mathrm{mod}\ 8),
\]
we set:
\[
n = \frac{m+1}{2} \quad \text{or} \quad \frac{m+9}{2},
\]
according as $m \equiv 1$ or $\equiv 9\ (\mathrm{mod}\ 16)$, so that $n \equiv 1\ (\mathrm{mod}\ 8)$ as well; then one has to set $x = 1$, $y = \pm 1$ or $= \pm 3$, and if we assume $A$ relatively prime to $n$, then $\partial = 1$, $a = n$, and from (18) it follows: % [sic] German print omits y before the second equality
\[
\begin{aligned}
\alpha + \beta &= A, & n(\gamma - \delta) &= -C + c(B+y), \\
\alpha - \beta &= B + y, & n(\gamma + \delta) &= -B + y + cA.
\end{aligned}
\]
Accordingly, since $B+y$ is odd and $c \equiv 0\ (\mathrm{mod}\ 16)$, condition (17) gives:
\[
A \equiv C \quad (\mathrm{mod}\ 8),
\]
%CONT>>
